\par\bigskip
\SourceLeaf{421}
\begin{center}
\pdfbookmark[2]{Gauss an Olbers, 26. November 1810}{pallas-brief-18101126}
\textsc{Gauss} an \textsc{Olbers}, 26. November 1810.
\end{center}
\noindent\hspace*{\parindent}Ich glaube Ihnen bereits in meinem letzten Briefe geschrieben zu haben, dass ich Willens war, eine Arbeit\HistoricalNoteEight{Vgl. S. 473 dieses Bandes.]} über die Störungen der Pallas durch den Jupiter zu unternehmen, nicht die allgemeine Theorie derselben, sondern nur ihren Betrag seit 1803 bis 1811 zu entwickeln. Die Berechnung dreier Systeme von elliptischen Elementen aus den Oppositionen von 1803, 1804, 1805, 1807 — 1804, 1805, 1807, 1808 — 1805, 1807, 1808, 1809 hatte die starke Einwirkung der Störungen erwiesen; urtheilen Sie selbst nach folgenden Proben:
\begin{center}
\begingroup\footnotesize\setlength{\tabcolsep}{3pt}\renewcommand{\arraystretch}{1.2}
\begin{tabular}{lccc}
Perihelium 1803\ .\,.\,.\,.\,.\,.\,.\,. & \(121^{0}\,3'11''\) & \(121^{0}\,5'22''\) & \(120^{0}58'\,5''\)\\
\(\varphi\)\ .\,.\,.\,.\,.\,.\,.\,.\,.\,.\,.\,.\,.\,.\,.\,.\,.\,.\,.\,.\,. & \(14^{0}10'59''\) & \(14^{0}10'\,4''\) & \(14^{0}\,9'37''\)\\
Tägl. mittl. Bewegung.\ .\,.\,. & \(770''{,}214\) & \(770''{,}447\) & \(770''{,}926.\)\\
\end{tabular}\endgroup
\end{center}
\par
\noindent\hspace*{\parindent}Ich wurde neugierig, ob diese sechs Oppositionen sich viel besser vereinigen lassen würden, wenn die Störungen vom Jupiter mit in Betracht gezogen würden. Ich gestehe, dass eine solche an sich nicht angenehme Arbeit so lange, als Eine Ellipse alle Beobachtungen darstellte, zu wenig Interesse für mich gehabt hatte, und dass selbst der Preis von 6000 Francs mich kaum dazu hätte anlocken können. Jetzt habe ich indess jene Arbeit muthig unternommen, und nachdem ich 4 Wochen den grössten Theil meiner disponiblen Zeit dazu angewandt habe, die Hauptsache abgethan, so dass ich, wenn nicht Hauptfehler in meiner Rechnung sind, bald die Früchte davon zu pflücken hoffen darf. Ich habe es am zweckmässigsten gehalten, die Elemente als veränderlich zu betrachten, und ihre täglichen Änderungen durch Jupiter von 50 zu 50 Tagen durch einen Zeitraum von mehr als 8 Jahren zu berechnen. Ich fing damit an, mir selbst Formeln dazu zu entwickeln, oder vielmehr die schon vor mehrern Jahren von mir entwickelten in die möglichst geschmeidige Gestalt zu bringen; denn die \textsc{Laplace}schen hielt ich für viel zu unbequem, und ich hätte bei Anwendung der letztern gewiss zwei- bis dreimal so viel Zeit nöthig gehabt als bei den meinigen. Die Jupitersörter berechnete ich nach den \textsc{Bouvard}schen Tafeln (mit Anwendung eines besondern Kunstgriffs\HistoricalNoteDouble{Vgl. S. 475 dieses Bandes.]}) \SourceLeaf{422}und die Pallasörter nach einem 4. System von Elementen, das an alle Oppositionen sich möglichst genau anschliesst. Jene 61 täglichen Änderungen der einzelnen sechs Elemente habe ich nun endlich heute vollendet, so dass theils noch die mechanische Integration (die nicht viel mehr als Addiren ist), theils die neue Berechnung der Grundelemente, die mit den so sich ergebenden Störungen verbunden die Oppositionen am besten darstellen, übrig bleibt. Diese Störungen sind viel beträchtlicher, als ich dachte, und besonders überwiegend stark im Jahr 1809, wo Jupiter mit der Pallas in heliocentrischer Conjunction war, daher ich glaube, dass meine Ephemeride in der nächsten Opposition sehr viel abweichen wird. \textsc{Harding} hat neulich die Pallas gesucht, weiss aber noch nicht gewiss, ob er sie gefunden hat; sollten Sie sie beobachten, so verbinden Sie mich sehr durch baldige Mittheilung. Die tägliche Bewegung wird jetzt mehr als \(2''\) grösser sein als 1803; der Winkel \(\varphi\), welcher beständig abgenommen hat, wohl \(13'\) kleiner, das Perihel ist auch wohl um \(20'\) oder \(30'\) zurückgegangen etc. Sollte die Übereinstimmung mit den sechs bisherigen Oppositionen und der 7. des nächsten Jahrs zu meiner Satisfaction ausfallen, so unternehme ich vielleicht auch die allgemeine Theorie, obwohl das ungeheuer viele mechanische Rechnen dabei sehr abschreckend ist. Die Störungen durch Jupiter sind so überwiegend, dass man die durch die andern Planeten vorerst noch bei Seite setzen kann.\quad.\quad.\quad.\quad.\quad.
\par
\noindent\hspace*{\parindent}Sollten Sie durch Ihren Wink in Paris eine nochmalige Verlängerung der Preisfrage veranlassen, so würde ich dadurch einen Grund mehr zur Unternehmung der vorhin erwähnten Arbeit haben. Das Gelingen meiner gegenwärtigen wird indess die erste Bedingung sein müssen. In vierzehn Tagen werde ich darüber im Klaren sein, wenigstens in Beziehung auf die bisherigen Beobachtungen; doppelt neugierig werde ich dann auf die demnächstigen sein, wo sich die bei der Conjunction erlittenen Störungen erst in ihrer ganzen Stärke zeigen werden.\quad.\quad.\quad.\quad.\quad.
\par
\noindent\hspace*{\parindent}P. S. Vor Schluss dieses Briefs habe ich noch die Integration der Störung der Neigung vollendet, die am wenigsten beträchtlich ist. Setze ich sie in der 1. Opposition
\begin{center}1803\quad zu\quad \(34^{0}37'47'',\)\end{center}
so ist sie
\SourceLeaf{423}
\begin{center}
\begin{tabular}{lc}
1804 & \(34^{0}37'36''\)\\
1805 & \(34^{0}37'12''\)\\
1807 & \(34^{0}37'25''\)\\
1808 & \(34^{0}37'31''\)\\
1809 & \(34^{0}36'40''\)\\
1811 & \(34^{0}35'\,5''.\)\\
\end{tabular}
\end{center}
\par
